\subsection{現在時点の経済的意思決定}

現在時点の経済的意思決定は、時間価値を扱わない判断である。将来の利率や割引ではなく、現在の費用関数や資源配分を使って判断する。

\subsubsection{損益分岐分析}

損益分岐分析は、複数案の費用が等しくなる点を見つける方法である。利用量が分岐点より小さい場合は案 A が安く、利用量が分岐点より大きい場合は案 B が安い、といった判断に使う。

\begin{notebox}{例}
クラウド事業者 A は月額固定費が低いが利用量単価が高い。事業者 B は月額固定費が高いが利用量単価が低い。この場合、利用量が少なければ A、多ければ B が有利になる。両者の費用が同じになる利用量が損益分岐点である。
\end{notebox}

\subsubsection{最適化分析}

最適化分析は、一つ以上の費用関数を調べ、総費用が最小になる点を見つける方法である。古典的な例は、空間と時間のトレードオフである。高速なアルゴリズムは多くのメモリを使うことがあり、メモリ費用と実行時間短縮の価値を比較する必要がある。

\par\smallskip
\begin{center}
\centering
\IfFileExists{assets/generated_figures/ch15/fig-ch15-05.png}{\includegraphics[width=0.96\linewidth]{assets/generated_figures/ch15/fig-ch15-05.png}}{\fbox{\parbox[c][48mm][c]{0.9\linewidth}{\centering 画像生成待ち\\損益分岐点と最適化の概念図}}}
\par\smallskip
\refstepcounter{figure}\label{fig:ch15-05}図 \thefigure: 損益分岐点と最適化の概念図
\end{center}
図 \thefigure は、損益分岐点と最適化の概念図を視覚的に整理したものである。
\par\smallskip
